\documentclass[aspectratio=169,12pt]{beamer}
% Build from the repository root: make lecture-06-slides
\usepackage[T1]{fontenc}
\usepackage{lmodern}
\usepackage{amsmath,amssymb,booktabs,array}
\usepackage{pgfplots}
\pgfplotsset{compat=1.18}
\usefonttheme{professionalfonts}
\definecolor{Ink}{HTML}{16324F}
\definecolor{Blue}{HTML}{1F4E79}
\definecolor{Orange}{HTML}{C46A13}
\definecolor{Teal}{HTML}{007D80}
\definecolor{Muted}{HTML}{526477}
\setbeamercolor{normal text}{fg=Ink,bg=white}
\setbeamercolor{frametitle}{fg=Ink}
\setbeamercolor{structure}{fg=Blue}
\setbeamercolor{footline}{fg=Muted}
\setbeamerfont{frametitle}{size=\large,series=\bfseries}
\setbeamerfont{footline}{size=\tiny}
\setbeamersize{text margin left=8mm,text margin right=8mm}
\setbeamertemplate{navigation symbols}{}
\setbeamertemplate{headline}{}
\setbeamertemplate{frametitle}{\vspace{3mm}\insertframetitle\par\vspace{1mm}}
\newcommand{\conspectsection}{6.1}
\setbeamertemplate{footline}{%
  \hspace{8mm}\usebeamerfont{footline}\usebeamercolor[fg]{footline}%
  Conspect \S\,\conspectsection\hfill\insertframenumber\hspace{8mm}\vspace{3mm}}
\setlength{\parskip}{0.5em}
\setlength{\abovedisplayskip}{0.55em}
\setlength{\belowdisplayskip}{0.55em}
\newcommand{\Prob}{\mathbb{P}}
\newcommand{\Expect}{\mathbb{E}}
\DeclareMathOperator{\Var}{Var}
\DeclareMathOperator{\Cov}{Cov}
\newcommand{\takeaway}[1]{\vfill{\small\color{Muted}#1\par}}
\newcommand{\plotcaption}[1]{\par\smallskip{\small #1\par}}
\setbeamertemplate{blocks}[default]
\setbeamerfont{block title}{size=\normalsize,series=\bfseries}
\setbeamerfont{block body}{size=\normalsize}
\newenvironment{coursebox}[2]{%
  \begingroup
  \setbeamercolor{block title}{fg=#1,bg=#1!13!white}
  \setbeamercolor{block body}{fg=Ink,bg=#1!4!white}
  \begin{block}{#2}
  \setlength{\parskip}{0.2em}
  \setlength{\abovedisplayskip}{0.3em}
  \setlength{\belowdisplayskip}{0.3em}
}{\end{block}\endgroup}
\newenvironment{definitionbox}[1]
  {\begin{coursebox}{Blue}{Definition: #1}}{\end{coursebox}}
\newenvironment{resultbox}[1]
  {\begin{coursebox}{Teal}{Result: #1}}{\end{coursebox}}
\newcommand{\InfDataPath}{../../figures/data}
\newcommand{\InfPlotWidth}{12cm}
\newcommand{\InfPlotHeight}{3.6cm}
% Shared Chapter 6 scientific charts: analytic curves and reproducible
% simulations. Regenerate the committed tables with generate_chapter6_figures.py.
% The conspect and slides may override these dimensions and the data path.
\providecommand{\InfDataPath}{figures/data}
\providecommand{\InfPlotWidth}{11.8cm}
\providecommand{\InfPlotHeight}{4.6cm}
\definecolor{InfBlue}{HTML}{1F4E79}
\definecolor{InfOrange}{HTML}{C46A13}
\definecolor{InfTeal}{HTML}{007D80}
\definecolor{InfMuted}{HTML}{526477}
\pgfplotsset{inf axis/.style={
  scale only axis,width=\InfPlotWidth,height=\InfPlotHeight,
  axis lines=left,enlargelimits=false,
  grid=major,grid style={black!9},
  tick label style={font=\footnotesize},label style={font=\footnotesize},
  legend style={draw=none,fill=none,font=\footnotesize,
    at={(0.5,1.05)},anchor=south,legend columns=-1,column sep=0.8em},
  every axis plot/.append style={no marks,line width=1.15pt},
}}

\newcommand{\InfCoveragePlot}{%
\begin{tikzpicture}\begin{axis}[inf axis,
  xmin=7.8,xmax=12.4,ymin=0,ymax=41,
  xlabel={Candidate population mean $\mu$},ylabel={Repeated dataset},
  xtick={8,9,10,11,12},ytick={1,10,20,30,40},
  ymajorgrids=false]
  \addplot[InfMuted,densely dashed,forget plot]
    coordinates {(10,0) (10,41)};
  \addplot[InfTeal,only marks,mark=*,mark size=0.65pt,
    error bars/.cd,x dir=both,x explicit,
    error bar style={line width=0.6pt},error mark options={rotate=90,mark size=1.0pt}]
    table[x=mean,y=trial,x error=halfwidth] {\InfDataPath/ch6-coverage-hit.dat};
  \addlegendentry{Contains $\mu=10$}
  \addplot[InfOrange,only marks,mark=*,mark size=0.85pt,
    error bars/.cd,x dir=both,x explicit,
    error bar style={line width=1.0pt},error mark options={rotate=90,mark size=1.2pt}]
    table[x=mean,y=trial,x error=halfwidth] {\InfDataPath/ch6-coverage-miss.dat};
  \addlegendentry{Misses $\mu=10$}
\end{axis}\end{tikzpicture}%
}

\newcommand{\InfStudentPlot}{%
\begin{tikzpicture}\begin{axis}[inf axis,
  xmin=-4,xmax=4,ymin=0,ymax=0.425,
  xlabel={Standardized estimation error},ylabel={Density},
  xtick={-4,-2,0,2,4},ytick={0,0.1,0.2,0.3,0.4}]
  \addplot[draw=none,fill=InfOrange!18,forget plot]
    table[x=x,y=lefttail] {\InfDataPath/ch6-student.dat} \closedcycle;
  \addplot[draw=none,fill=InfOrange!18,forget plot]
    table[x=x,y=righttail] {\InfDataPath/ch6-student.dat} \closedcycle;
  \addplot[InfBlue] table[x=x,y=normal] {\InfDataPath/ch6-student.dat};
  \addlegendentry{Standard Normal}
  \addplot[InfOrange,dashed] table[x=x,y=t9] {\InfDataPath/ch6-student.dat};
  \addlegendentry{Student $t_9$}
  \addplot[InfBlue,densely dotted,forget plot]
    coordinates {(-1.959963985,0) (-1.959963985,0.05844507)};
  \addplot[InfBlue,densely dotted,forget plot]
    coordinates {(1.959963985,0) (1.959963985,0.05844507)};
  \addplot[InfOrange,densely dotted,forget plot]
    coordinates {(-2.262157163,0) (-2.262157163,0.0407427)};
  \addplot[InfOrange,densely dotted,forget plot]
    coordinates {(2.262157163,0) (2.262157163,0.0407427)};
  \node[anchor=north east,font=\footnotesize,text=InfMuted,align=right,
    fill=white,inner sep=2pt] at (axis description cs:0.98,0.97)
    {$z_{0.975}=1.960$\\$t_{9,0.975}=2.262$};
\end{axis}\end{tikzpicture}%
}

\newcommand{\InfProportionCoveragePlot}{%
\begin{tikzpicture}\begin{axis}[inf axis,
  xmin=0,xmax=1,ymin=0,ymax=1.03,
  xlabel={True Bernoulli probability $p$},ylabel={Exact coverage probability},
  xtick={0,0.2,0.4,0.6,0.8,1},ytick={0,0.25,0.5,0.75,0.95}]
  \addplot[InfMuted,densely dotted]
    coordinates {(0,0.95) (1,0.95)};
  \addlegendentry{Nominal $95\%$}
  \addplot[InfOrange,line width=0.85pt]
    table[x=p,y=wald] {\InfDataPath/ch6-proportion-coverage.dat};
  \addlegendentry{Wald, $n=20$}
  \addplot[InfTeal,line width=1.0pt]
    table[x=p,y=wilson] {\InfDataPath/ch6-proportion-coverage.dat};
  \addlegendentry{Wilson, $n=20$}
\end{axis}\end{tikzpicture}%
}

\newcommand{\InfBootstrapPlot}{%
\begin{tikzpicture}\begin{axis}[inf axis,
  xmin=2.3,xmax=7,ymin=0,ymax=0.75,
  xlabel={Bootstrap sample mean $\bar X^*$},ylabel={Bootstrap density},
  xtick={2.5,3.4,4.5,5.68,7},
  xticklabels={$2.5$,$3.40$,$4.5$,$5.68$,$7$}]
  \addplot[InfBlue,fill=InfBlue!20,const plot,line width=0.65pt]
    table[x=value,y=density] {\InfDataPath/ch6-bootstrap-mean.dat} \closedcycle;
  \addlegendentry{$20\,000$ resamples of $n=40$}
  \addplot[InfOrange,dashed,forget plot]
    coordinates {(3.399081569,0) (3.399081569,0.7)};
  \addplot[InfOrange,dashed,forget plot]
    coordinates {(5.681750915,0) (5.681750915,0.7)};
  \addplot[InfTeal,densely dotted,forget plot]
    coordinates {(4.494628920,0) (4.494628920,0.7)};
  \node[anchor=north east,font=\footnotesize,text=InfMuted,align=right,
    fill=white,inner sep=2pt] at (axis description cs:0.98,0.98)
    {Observed mean: $4.495$\\Percentile CI: $[3.40,5.68]$};
\end{axis}\end{tikzpicture}%
}

\newcommand{\InfPairedPlot}{%
\begin{tikzpicture}\begin{axis}[inf axis,
  xmin=-0.02,xmax=0.23,ymin=0,ymax=0.072,
  xlabel={Bootstrap accuracy difference $\widehat a_A^*-\widehat a_B^*$},
  ylabel={Bootstrap probability},xtick={0,0.04,0.1,0.16,0.2},
  xticklabels={$0$,$0.04$,$0.10$,$0.16$,$0.20$},
  ytick={0,0.02,0.04,0.06},scaled y ticks=false,
  yticklabel={\pgfmathprintnumber[fixed,precision=2]{\tick}}]
  \addplot[InfBlue,ycomb,line width=1.7pt]
    table[x=difference,y=mass] {\InfDataPath/ch6-paired-accuracy.dat};
  \addlegendentry{Paired bootstrap, $n=200$, $B=20\,000$}
  \addplot[InfMuted,densely dotted,forget plot]
    coordinates {(0,0) (0,0.068)};
  \addplot[InfOrange,dashed,forget plot]
    coordinates {(0.04,0) (0.04,0.068)};
  \addplot[InfOrange,dashed,forget plot]
    coordinates {(0.16,0) (0.16,0.068)};
  \node[anchor=north east,font=\footnotesize,text=InfMuted,align=right,
    fill=white,inner sep=2pt] at (axis description cs:0.98,0.98)
    {Observed gap: $0.10$\\Percentile CI: $[0.04,0.16]$};
\end{axis}\end{tikzpicture}%
}

\title{Sampling distributions and confidence intervals}
\subtitle{Lecture 6: Bootstrap and uncertainty after fitting}
\author{}
\date{Scientific Computing 2026}
\hypersetup{pdftitle={Lecture 6: Sampling distributions, confidence intervals, and bootstrap},
  pdfsubject={Scientific Computing 2026, Chapter 6},pdfauthor={}}

\begin{document}

% 1
{\setbeamercolor{background canvas}{bg=Ink}
\begin{frame}[plain]
\color{white}
{\large Scientific Computing 2026\par}
\vfill
{\LARGE\bfseries Sampling distributions,\par
confidence intervals, and bootstrap\par}
\vspace{0.7em}
{\large Uncertainty after fitting\par}
\vfill
{\large Lecture 6\par}\vspace{4mm}
\end{frame}}

% 2
\begin{frame}{Five waiting times, one fitted rate}
The observations from Chapter 5 were
\[
  1,\ 2,\ 3,\ 4,\ 5\text{ minutes}.
\]
The exponential model gave
\[
  \widehat\lambda=\frac1{\bar t}=\frac13\text{ per minute}.
\]
How much would that estimate change if we collected five new waits?
\end{frame}

% 3
\begin{frame}{The parameter stays fixed when the sample changes}
Imagine repeating the observation process under one true rate $\lambda$.
\[
  (T_1,\ldots,T_n)\quad\longmapsto\quad
  \widehat\lambda=\frac n{\sum_iT_i}.
\]
Each new sample produces a new estimate. Their distribution describes
the estimation error under that rate.
\vfill
In practice, we have one dataset and do not know $\lambda$.
\end{frame}

% 4
\begin{frame}{The quantity whose uncertainty matters}
An interval must refer to a specific population quantity.
\par\smallskip
{\renewcommand{\arraystretch}{1.5}
\begin{tabular*}{\textwidth}{@{\extracolsep{\fill}}ll@{}}
\toprule
Data & Possible target\\
\midrule
Waiting times & Arrival rate $\lambda$ or mean wait $1/\lambda$\\
Repeated measurements & Population mean $\mu$\\
Independent test examples & Accuracy $p$ of a fixed classifier\\
\bottomrule
\end{tabular*}}\par
\takeaway{The variation of individual outcomes and the uncertainty of an
estimated population quantity answer different questions.}
\end{frame}

\renewcommand{\conspectsection}{6.2}
% 5
\begin{frame}{Confidence intervals}
\begin{definitionbox}{Confidence interval with level $1-\alpha$}
A rule that takes a sample $X=(X_1,\ldots,X_n)$ and returns
$[L(X),U(X)]$, satisfying
\[
  \Prob_\theta\{L(X)\le\theta\le U(X)\}\ge1-\alpha
  \qquad\text{for every allowed }\theta.
\]
The probability concerns repeated samples from $P_\theta$.
\end{definitionbox}
An \textbf{exact} guarantee holds at the stated sample size under the model.
An \textbf{approximate} interval can have coverage different from its nominal level.
\end{frame}

% 6
\begin{frame}{Ninety-five percent coverage}
\centering\InfCoveragePlot
\plotcaption{Each of forty Gaussian samples ($n=25$, known $\sigma=2$)
gives an interval. The fixed true mean is $10$. Orange intervals miss it.}
\end{frame}

% 7
\begin{frame}{What the confidence level means}
Before seeing the data, both endpoints are random.
The interval construction covers the fixed parameter in at least $95\%$
of repeated samples when it has a $95\%$ guarantee.
\vfill
After observing the data, the endpoints have numerical values.
The realized interval either contains the parameter or misses it.
\takeaway{A $95\%$ confidence level does not assign a posterior probability
of $95\%$ to the parameter lying in the realized interval.}
\end{frame}

% 8
\begin{frame}{Quantiles choose the central probability}
\begin{definitionbox}{Quantile}
For a CDF $F$ and $0<p<1$,
\[
  Q_p(F)=\inf\{x:F(x)\ge p\}.
\]
For a continuous, strictly increasing CDF, $F(Q_p(F))=p$.
\end{definitionbox}
Write $z_p$ for the standard Gaussian quantile.
\[
  z_{0.975}\approx1.96,\qquad
  \Prob(-1.96\le Z\le1.96)\approx0.95.
\]
\end{frame}

% 9
\begin{frame}{A Gaussian mean with known variance}
Assume $X_i\overset{\mathrm{iid}}\sim\mathcal N(\mu,\sigma^2)$,
with known $\sigma>0$.
\[
  Z=\frac{\sqrt n(\bar X-\mu)}{\sigma}\sim\mathcal N(0,1).
\]
The standardized error has a known law that does not depend on $\mu$.
Such a quantity is called a \textbf{pivot}.
\[
  \Prob_\mu\!\left\{-z_{1-\alpha/2}\le
  \frac{\sqrt n(\bar X-\mu)}\sigma
  \le z_{1-\alpha/2}\right\}=1-\alpha.
\]
\end{frame}

% 10
\begin{frame}{An interval by rearranging the inequalities}
\begin{resultbox}{Gaussian mean, known variance}
Under the preceding Gaussian assumptions, an exact interval is
\[
  \left[\bar X-z_{1-\alpha/2}\frac\sigma{\sqrt n},\quad
        \bar X+z_{1-\alpha/2}\frac\sigma{\sqrt n}\right].
\]
\end{resultbox}
Larger $\sigma$ produces a wider interval.\
Four times as many observations halves its width.\
Higher confidence increases the quantile and widens the interval.
\end{frame}

\renewcommand{\conspectsection}{6.3}
% 11
\begin{frame}{The noise variance usually needs estimation}
Use the sample variance from Chapter 5,
\[
  S^2=\frac1{n-1}\sum_i(X_i-\bar X)^2,\qquad n\ge2.
\]
Replacing $\sigma$ by $S$ also puts random variation into the denominator:
\[
  \frac{\sqrt n(\bar X-\mu)}S.
\]
An unusually small $S$ can make this standardized error large.
The Gaussian cutoff $1.96$ no longer gives the exact $95\%$ interval.
\end{frame}

% 12
\begin{frame}{The distribution of a Gaussian sum of squares}
\begin{definitionbox}{Chi-square distribution}
For independent $Z_j\sim\mathcal N(0,1)$,
\[
  U=\sum_{j=1}^{\nu}Z_j^2\sim\chi^2_\nu.
\]
This is Gamma with shape $\nu/2$ and rate $1/2$.
The number $\nu$ is called the degrees of freedom.
\end{definitionbox}
For an i.i.d. Gaussian sample,
\[
  U=\frac{(n-1)S^2}{\sigma^2}\sim\chi^2_{n-1},
  \qquad U\ \text{is independent of }\bar X.
\]
\takeaway{Estimating the sample mean leaves $n-1$ independent directions
for the centered Gaussian variation.}
\end{frame}

% 13
\begin{frame}{The Student pivot for an unknown variance}
The Student construction from Chapter 2 takes the form
\[
  \frac{Z}{\sqrt{U/\nu}}\sim t_\nu,
  \quad Z\sim\mathcal N(0,1),\quad U\sim\chi^2_\nu,
  \quad Z\perp U.
\]
Here $\nu=n-1$, so
\[
  \frac{\sqrt n(\bar X-\mu)}S\sim t_{n-1}.
\]
Write $t_{\nu,p}$ for the $p$-quantile of $t_\nu$.
\begin{resultbox}{Gaussian mean, unknown variance}
For i.i.d. Gaussian data, an exact $1-\alpha$ interval is
\[
  \bar X\ \pm\ t_{n-1,\,1-\alpha/2}\frac S{\sqrt n}
\]
\end{resultbox}
\end{frame}

% 14
\begin{frame}{Estimating variance makes the tails heavier}
\centering\InfStudentPlot
\plotcaption{The Student law gives more room to large standardized errors.
For nine degrees of freedom, the $97.5\%$ quantile is approximately $2.262$.}
\end{frame}

% 15
\begin{frame}{Ten measurements}
Suppose the Gaussian model is appropriate and
\[
  n=10,\qquad\bar x=10.2,\qquad s=1.5.
\]
The $95\%$ Student interval has margin
\[
  t_{9,0.975}\frac s{\sqrt n}
  \approx2.262\frac{1.5}{\sqrt{10}}\approx1.073.
\]
\begin{resultbox}{95\% confidence interval for the population mean}
\[
  C(x)=[9.13,\ 11.27].
\]
\end{resultbox}
\takeaway{Using $1.96$ instead gives a smaller margin of $0.930$,
which ignores the uncertainty in estimating the variance.}
\end{frame}

% 16
\begin{frame}{Mean intervals beyond Gaussian data}
For i.i.d. observations with finite, positive variance,
the studentized CLT gives
\[
  \frac{\sqrt n(\bar X-\mu)}S
  \ \xrightarrow{d}\ \mathcal N(0,1).
\]
Thus $\bar X\pm1.96S/\sqrt n$ is an asymptotic $95\%$ interval.
The same Student formula often gives a similar approximation for large $n$.
\vfill
For a finite non-Gaussian sample, the exact $t_{n-1}$ law need not hold.
\takeaway{Skewness and heavy tails can make the approximation poor.
There is no universal sample size such as $n=30$ that guarantees accuracy.}
\end{frame}

% 17
\begin{frame}{A population mean and a future measurement}
The mean interval becomes narrow as we estimate $\mu$ more precisely:
\[
  \bar X\pm t_{n-1,\,1-\alpha/2}\frac S{\sqrt n}.
\]
A future observation also has its own measurement noise.
Under the same i.i.d. Gaussian model, an exact prediction interval is
\[
  \bar X\pm t_{n-1,\,1-\alpha/2}S\sqrt{1+\frac1n}.
\]
\takeaway{The prediction interval remains wide even with a very large
sample. Individual measurements still vary around the mean.}
\end{frame}

% 18
\begin{frame}{An exact interval for the exponential rate}
For independent exponential waits, their random total satisfies
\[
  W=\sum_iT_i\sim\operatorname{Gamma}(n,\lambda),
  \qquad2\lambda W\sim\chi^2_{2n}.
\]
Write $c_{\nu,p}$ for the $p$-quantile of $\chi^2_\nu$.
\begin{resultbox}{Exponential rate interval}
\[
  \left[\frac{c_{2n,\alpha/2}}{2W},\quad
        \frac{c_{2n,1-\alpha/2}}{2W}\right].
\]
\end{resultbox}
For the original five waits, the observed total is $w=15$ minutes:
\[
  \widehat\lambda=0.333,\qquad
  95\%\text{ interval }[0.108,\ 0.683]\text{ per minute}.
\]
\end{frame}

\renewcommand{\conspectsection}{6.4}
% 19
\begin{frame}{Accuracy on independent test examples}
For a fixed classifier, let $Y_i=1$ when its prediction is correct.
Under independent examples from one test population,
\[
  Y_i\sim\operatorname{Bernoulli}(p),\qquad
  \widehat p=\frac kn,\qquad
  \operatorname{SE}(\widehat p)=\sqrt{\frac{p(1-p)}n}.
\]
The CLT and substitution of $\widehat p$ give the Wald interval:
\[
  \widehat p\ \pm\ z_{1-\alpha/2}
  \sqrt{\frac{\widehat p(1-\widehat p)}n}.
\]
\takeaway{For $k=84$ correct predictions out of $n=100$,
the approximate $95\%$ interval is $[0.768,\ 0.912]$.}
\end{frame}

% 20
\begin{frame}{Why the Wald interval can fail}
Suppose zero out of twenty tested cases have the outcome of interest.
Then $\widehat p=0$, and the estimated standard error is also zero.
\[
  \text{Wald interval}=[0,0].
\]
But a small positive probability can easily produce twenty zero outcomes:
\[
  p=0.05\quad\Longrightarrow\quad
  \Prob(k=0)=0.95^{20}\approx0.358.
\]
\takeaway{Substituting a boundary estimate into the variance collapses
the interval. Small counts also make Gaussian approximations unreliable.}
\end{frame}

% 21
\begin{frame}{An interval that retains the candidate variance}
For each candidate probability $p$, use its variance in the standardized error:
\[
  \frac{\widehat p-p}{\sqrt{p(1-p)/n}}.
\]
Keep the candidates for which the error is at most $z=z_{1-\alpha/2}$ in magnitude:
\[
  n(\widehat p-p)^2\le z^2p(1-p).
\]
Solving this quadratic inequality for $p$ gives the Wilson interval.
\takeaway{This construction also uses a Gaussian approximation.
Its coverage is generally much better behaved near the boundaries.}
\end{frame}

% 22
\begin{frame}{The Wilson proportion interval}
\begin{resultbox}{Wilson interval, nominal level $1-\alpha$}
With $z=z_{1-\alpha/2}$, the interval is
\[
  \frac{\displaystyle\widehat p+\frac{z^2}{2n}
    \ \pm\ z\sqrt{\frac{\widehat p(1-\widehat p)}n+
                      \frac{z^2}{4n^2}}}
       {\displaystyle1+\frac{z^2}n}.
\]
\end{resultbox}
For zero outcomes in twenty trials, the $95\%$ interval is
\[
  \left[0,\ \frac{1.96^2}{20+1.96^2}\right]
  \approx[0,\ 0.161].
\]
\end{frame}

% 23
\begin{frame}{Nominal confidence and actual coverage}
\centering\InfProportionCoveragePlot
\plotcaption{For $n=20$ and each true $p$, sum binomial probabilities of counts
whose interval contains $p$. Discrete counts make coverage fluctuate.}
\end{frame}

\renewcommand{\conspectsection}{6.5}
% 24
\begin{frame}{Using the observed distribution}
When a statistic has no convenient sampling formula, the observations
can supply an approximate population law.
\begin{definitionbox}{Empirical distribution}
Put probability $1/n$ on each observed row:
\[
  F_n(t)=\frac1n\sum_{i=1}^n\mathbf1_{\{x_i\le t\}}.
\]
Repeated values receive the sum of their masses.
\end{definitionbox}
Drawing from this law means selecting an observed row uniformly at random.
\end{frame}

% 25
\begin{frame}{One bootstrap sample}
Start with the five observed waits
\[
  x=(1,2,3,4,5),\qquad\bar x=3.
\]
Draw five independent indices uniformly from $\{1,\ldots,5\}$,
allowing repeated indices. One possible result is
\[
  x^*=(1,1,2,4,5),\qquad\bar x^*=2.6.
\]
Other draws omit some observations and repeat others.
\takeaway{The star denotes a random resample, conditional on the fixed
observed data. Each bootstrap sample has the original sample size $n$.}
\end{frame}

% 26
\begin{frame}{Bootstrap replicates of a fitted statistic}
For a statistic $\widehat\theta=T(x_1,\ldots,x_n)$:
\begin{enumerate}
  \item Draw $n$ indices independently and uniformly, with replacement.
  \item Form the resampled rows $x_1^*,\ldots,x_n^*$.
  \item Recompute the whole statistic, $\widehat\theta^*=T(x_1^*,\ldots,x_n^*)$.
  \item Repeat to get $B$ replicates $\widehat\theta^{*(1)},\ldots,
        \widehat\theta^{*(B)}$.
\end{enumerate}
\takeaway{The same fitting rule applies to every replicate.
The examples here use $B=20{,}000$ to stabilize the simulated quantiles.}
\end{frame}

% 27
\begin{frame}{The error distribution we try to approximate}
The real estimation error is $\widehat\theta-\theta$.
In the bootstrap population, the fitted counterpart of $\theta$ is $\widehat\theta$.
The bootstrap error is therefore
\[
  \widehat\theta^*-\widehat\theta.
\]
For an i.i.d. sample mean with finite, positive population variance,
the centered bootstrap error approximates the sampling error as $n$ grows:
\[
  \mathcal L\!\left(\sqrt n(\bar X^*-\bar x)\mid x\right)
  \ \approx\ \mathcal L_\mu\!\left(\sqrt n(\bar X-\mu)\right).
\]
Here $\mathcal L$ denotes a distribution; conditioning on $x$ fixes the
observed data. Other statistics require their own approximation conditions.
\end{frame}

% 28
\begin{frame}{A bootstrap standard error}
\begin{resultbox}{Estimated spread of the statistic}
With $\overline\theta^*=B^{-1}\sum_b\widehat\theta^{*(b)}$,
\[
  \widehat{\operatorname{SE}}_{\mathrm{boot}}
  =\sqrt{\frac1{B-1}\sum_{b=1}^B
          (\widehat\theta^{*(b)}-\overline\theta^*)^2}.
\]
\end{resultbox}
For the five waits, the exact conditional bootstrap calculation gives
\[
  \Expect^*[\bar X^*]=3,\qquad
  \Var^*(\bar X^*)=\frac25=0.4.
\]
\takeaway{The usual sample estimate $S^2/n$ is $0.5$ here.
The empirical distribution uses denominator $n$ in its own variance.}
\end{frame}

% 29
\begin{frame}{Percentile bootstrap intervals}
Let $q_p^*$ denote the empirical $p$-quantile of the bootstrap estimates.
\begin{definitionbox}{Percentile interval}
\[
  [q_{\alpha/2}^*,\ q_{1-\alpha/2}^*].
\]
Keep the central $1-\alpha$ fraction of the bootstrap estimates.
\end{definitionbox}
This is a direct interval construction. Its coverage is approximate
and depends on the statistic and sampling distribution.
\end{frame}

% 30
\begin{frame}{Inverting centered bootstrap errors}
The approximate central error bounds are
\[
  q_{\alpha/2}^*-\widehat\theta
  \ \lesssim\ \widehat\theta-\theta
  \ \lesssim\ q_{1-\alpha/2}^*-\widehat\theta.
\]
Rearranging for $\theta$ gives another construction.
\begin{definitionbox}{Basic bootstrap interval}
\[
  [2\widehat\theta-q_{1-\alpha/2}^*,\quad
   2\widehat\theta-q_{\alpha/2}^*].
\]
\end{definitionbox}
\takeaway{The basic interval reflects the bootstrap quantiles about
$\widehat\theta$. It can differ from the percentile interval.}
\end{frame}

% 31
\begin{frame}{Bootstrap estimates of a mean waiting time}
\centering\InfBootstrapPlot
\plotcaption{Forty illustrative exponential waits have mean $4.495$ minutes.
The 20,000 bootstrap means give percentile interval $[3.399,\ 5.682]$ minutes.}
\end{frame}

\renewcommand{\conspectsection}{6.6}
% 32
\begin{frame}{Two classifiers on the same test examples}
Both fixed classifiers use the same 200 independent examples.
\par\smallskip
{\renewcommand{\arraystretch}{1.5}
\begin{tabular*}{\textwidth}{@{\extracolsep{\fill}}lrr@{}}
\toprule
 & B correct & B incorrect\\
\midrule
A correct & 150 & 30\\
A incorrect & 10 & 10\\
\bottomrule
\end{tabular*}}\par
\[
  \widehat p_A=0.90,\qquad\widehat p_B=0.80,\qquad
  \widehat\Delta=0.10.
\]
\takeaway{The estimated advantage is ten percentage points.
Its uncertainty depends on how the models' successes overlap.}
\end{frame}

% 33
\begin{frame}{A difference for each example}
Let $d_i=\mathbf1_{\{A\text{ correct}\}}-\mathbf1_{\{B\text{ correct}\}}$.
The rows give
\[
  d_i=\begin{cases}
    +1 & \text{on 30 examples},\\
    -1 & \text{on 10 examples},\\
     0 & \text{on 160 examples}.
  \end{cases}
\]
\[
  \widehat\Delta=\bar d=0.10,\qquad
  \widehat{\operatorname{SE}}(\bar d)
  =\sqrt{\frac{0.20-0.10^2}{199}}\approx0.0309.
\]
\takeaway{A Gaussian approximation gives $0.10\pm1.96(0.0309)$,
or approximately $[0.039,\ 0.161]$.}
\end{frame}

% 34
\begin{frame}{Bootstrap the paired test rows}
Each row contains one example and both classifiers' results:
\[
  (y_i,\widehat y_{A,i},\widehat y_{B,i}).
\]
Draw one list of row indices and use it for both models.
In each bootstrap sample, recompute
\[
  \widehat\Delta^*=\widehat p_A^*-\widehat p_B^*.
\]
Resampling the models independently destroys their observed association.
\takeaway{Indeed, $\Var(\widehat p_A-\widehat p_B)
=\Var(\widehat p_A)+\Var(\widehat p_B)
-2\Cov(\widehat p_A,\widehat p_B)$.}
\end{frame}

% 35
\begin{frame}{Paired bootstrap differences}
\centering\InfPairedPlot
\plotcaption{20,000 paired resamples give percentile interval $[0.04,\ 0.16]$.
Both fitted models stay fixed. Only the sampled test examples change.}
\end{frame}

% 36
\begin{frame}{Metrics built from several counts}
Let $TP$, $FP$, and $FN$ count true positives, false positives,
and false negatives. Then
\[
  \operatorname{precision}=\frac{TP}{TP+FP},\qquad
  \operatorname{recall}=\frac{TP}{TP+FN},\qquad
  F_1=\frac{2TP}{2TP+FP+FN}.
\]
Resample complete test rows and recompute all counts, then the metric.
For a model comparison, recompute both metrics using the same rows.
\vfill
If a denominator is zero, specify what the metric means before resampling.
Rare positive cases can make the bootstrap distribution unstable.
\takeaway{Averaging per-row precision or per-row $F_1$ values does not
reproduce these ratios of aggregate counts.}
\end{frame}

\renewcommand{\conspectsection}{6.7}
% 37
\begin{frame}{More resamples and more observations}
The original sample size $n$ controls how much information the data contain.
The number of bootstrap replicates $B$ controls simulation accuracy.
\[
  \text{Larger }B\quad\Longrightarrow\quad
  \text{more stable simulated standard errors and quantiles}.
\]
Once those summaries stabilize, further resampling leaves the statistical
uncertainty essentially unchanged.
\end{frame}

% 38
\begin{frame}{The resampling unit follows the observation process}
Ordinary row resampling assumes that rows act as independent sampling units.
\vfill
Repeated measurements from one user share user-specific effects.
Resampling users together with their measurements preserves that grouping.
\vfill
Nearby time points can be dependent. Resampling suitable consecutive blocks
can preserve some of that dependence.
\takeaway{Independent row resampling can underestimate uncertainty for dependent data.}
\end{frame}

% 39
\begin{frame}{Observed data can miss possible outcomes}
With zero recorded rare events, every ordinary bootstrap sample also has zero.
The empirical distribution cannot generate an unseen event type.
\vfill
For i.i.d. $X_i\sim\operatorname{Uniform}(0,\theta)$, let $M=\max_iX_i$.
Then $M<\theta$ almost surely. Every resampled maximum satisfies
\[
  M^*\le M,\qquad\Prob^*(M^*\le M)=1.
\]
The upper percentile endpoint is at most $M$, so the percentile interval
for $\theta$ has coverage zero.
\end{frame}

% 40
\begin{frame}{Test uncertainty and training uncertainty}
Resampling a test set while holding the model fixed measures uncertainty
from which test examples we observed.
\vfill
To study variation in the fitted model, resample training data and refit the
estimator in every replicate.
\vfill
If model selection is part of the procedure being studied, reproduce that
selection too. For the fixed-model accuracy claim in Chapter 4, evaluate
on independent test observations that were not used to fit or select the model.
\end{frame}

% 41
\begin{frame}{The sample must match the population claim}
If the dataset covers only one type of user, resampling it cannot reveal
performance for absent user groups.
\vfill
Sampling bias and distribution shifts require information or assumptions
about the missing population. Repeated use of the same observed rows
does not supply that information.
\end{frame}

% 42
\begin{frame}{An interval statement includes its assumptions}
For the ten Gaussian measurements,
\[
  \widehat\mu=10.2,\qquad95\%\text{ confidence interval }[9.13,11.27].
\]
This reports uncertainty about the \textbf{population mean}, assuming
independent measurements from one Gaussian distribution.
\vfill
The Student construction gives exact coverage under that model.
A bootstrap or CLT interval needs its corresponding approximation conditions.
\end{frame}

\end{document}
